\section{Main Theorems}

This section states the five core results. Complete proofs, equality/strictness conditions, and supplementary derivations are integrated in the manuscript appendix.

\subsection{Covariance identity}

For $X$ satisfying
\[
\mathbb E[|X|]<\infty,
\qquad
\mathbb E[|X|S]<\infty,
\]
and nonnegative $S$ with $0<\mathbb E[S]<\infty$,
\[
\mathbb E_{\mathrm{FP}}[X]-\mathbb E[X]
=
\frac{\operatorname{Cov}(X,S)}{\mathbb E[S]}.
\]
Hence positive first-person mean uplift is equivalent to positive outcome--accessibility covariance. The separate base and weighted integrability conditions are both needed in general.

\subsection{Tail identity}

For threshold $c$,
\[
P_{\mathrm{FP}}(X\ge c)-P(X\ge c)
=
\frac{\operatorname{Cov}(\mathbf 1_{\{X\ge c\}},S)}{\mathbb E[S]}.
\]

\subsection{Monotone accessibility and FOSD}

Let $g(x)=\mathbb E[S\mid X=x]$. If $g$ is nondecreasing, then
\[
F_{\mathrm{FP}}(c)\le F(c)
\]
for every $c$. Positive Pearson correlation alone is not sufficient; nonmonotone accessibility can produce CDF crossings.

\begin{figure}[t]
\centering
\includegraphics[width=0.82\linewidth]{../figures/generated_pdf/fig2_fosd.pdf}
\caption{Deterministic theorem illustration for the FOSD boundary. A monotone outcome-aligned accessibility rule shifts the first-person CDF weakly below the base CDF, whereas a nonmonotone rule can produce a crossing. The curves illustrate the theorem and its boundary; they are not empirical Everett data.}
\label{fig:fosd-boundary}
\end{figure}

\subsection{Recognition decomposition}

Under the formal-model assumptions
\[
0<\mathbb E[S_R]<\infty,
\qquad
\mathbb E[|U_R|]<\infty,
\qquad
\mathbb E[|U_R|S_R]<\infty
\]
for $R\in\{0,1\}$, recognition satisfies
\[
V_1-V_0
=
\mathbb E[U_1-U_0]
+
Q(U_1,S_1)-Q(U_0,S_0).
\]
If $S_0\equiv1$,
\[
V_1-V_0
=
\mathbb E[U_1-U_0]
+
\frac{\operatorname{Cov}(U_1,S_1)}{\mathbb E[S_1]}.
\]

\subsection{Policy--QBS interaction}

For the general selector-changing identity, also assume
\[
\mathbb E[|U_1|S_0]<\infty,
\]
so that the intermediate quantity $Q(U_1,S_0)$ is defined. Let $D=U_1-U_0$ and define
\[
I=Q(U_1,S_1)-Q(U_0,S_0).
\]
Then
\[
I
=
\frac{\operatorname{Cov}(D,S_0)}{\mathbb E[S_0]}
+
\left[Q(U_1,S_1)-Q(U_1,S_0)\right].
\]
For a fixed selector, the second term vanishes. If policy improvement increases with branch badness while accessibility decreases with branch badness, the interaction is nonpositive.

The core five remain fixed in identity and proof structure; current `main` includes explicit domain assumptions that close the audited integrability gaps in the frozen v0.3 wording. Two supplementary results are used for mechanism interpretation: S1 gives cross-branch policy coherence under explicit shared-latent assumptions, and S2 gives predictive-calibration alignment. In S2, if $m(Y)=\mathbb E[U\mid Y]$ and score-measurable accessibility $S=s(Y)$ are comonotone, then
\[
\operatorname{Cov}(U,S)
=
\operatorname{Cov}(m(Y),s(Y))
\ge0.
\]
The proof and the counterexample showing that positive mutual information alone is insufficient are kept in the supplementary appendix.
